\documentclass[10pt]{article}
\usepackage[a4paper,margin=20mm]{geometry}
\usepackage{amsmath,amssymb,amsthm,mathtools,bm}
\usepackage{mathrsfs}
\usepackage{booktabs,array,microtype,xurl,hyperref}
\hypersetup{colorlinks=true,linkcolor=blue,citecolor=blue,urlcolor=blue}

\newtheorem{theorem}{Theorem}[section]
\newtheorem{proposition}[theorem]{Proposition}
\newtheorem{lemma}[theorem]{Lemma}
\newtheorem{corollary}[theorem]{Corollary}
\newtheorem{definition}[theorem]{Definition}
\newtheorem{remark}[theorem]{Remark}
\newcommand{\dd}{\,\mathrm d}
\newcommand{\RR}{\mathbb R}
\newcommand{\SSS}{\mathbb S}
\newcommand{\cE}{\mathcal E}
\newcommand{\cA}{\mathcal A}
\newcommand{\cP}{\mathcal P}

\title{AP-E18 Graph-Tight Annuli on a Minimizing Diagonal:\\
Failure of the Reverse-H\"older Import, a Boundary-Straddling\\
Rank-Two Needle, and the Scale-Polished Repair}
\author{Route-F derivation ledger}
\date{24 July 2026}

\begin{document}
\maketitle

\begin{abstract}
AP-E17 reduced the volumetric relaxation gate to the graph-tight annulus
condition (GTA).  This note tests whether GTA follows from the actual
minimizing diagonal, rather than from an arbitrary bounded graph sequence.
The answer separates two notions of minimality.

The AP-E17 diagonal has an \emph{additive} comparison error bounded by its
global energy excess.  After normalization on a tangent ball this error is
$o(1)$, but it is not controlled relative to every unresolved microscopic
subball.  Consequently it gives no uniform reverse-H\"older seed.  Standard
$(p,q)$-growth regularity cannot be imported: in dimension three the current
ambient bounds have $(p,q)=(2,6)$, outside the 2024 relaxed-quasiconvex range
$q<\min\{3,3\}=3$.  Treating the complete-minor vector as an
$\mathscr A$-free quadratic field also fails, because the admissible graph is
cut out by nonlinear Pl\"ucker identities and the $\SSS^3$ constraint; the
map is not minimizing over the full linear $\mathscr A$-free class.

More decisively, we construct a smooth, boundary-straddling rank-two needle.
In tangent variables its primary amplitude, secondary amplitude, and radius
are
\[
 A_\varepsilon=\varepsilon^{-1/2},\qquad
 b_\varepsilon=h_\varepsilon=\varepsilon^2 .
\]
It has strong $L^2$ mass $O(\varepsilon^5)$, first- and second-minor energy
$O(\varepsilon)$, and zero pure third minor.  Nevertheless
\[
 \int |z_\varepsilon|^2|Dz_\varepsilon|^2\asymp1,\qquad
 \int |z_\varepsilon|^2|M_2(Dz_\varepsilon)|^2\asymp1 .
\]
Placed across the boundary of the tangent ball, it defeats \emph{every}
shrinking annulus thickness.  At physical tangent radius
$r_\varepsilon=\varepsilon$, its target amplitude is only
$r_\varepsilon A_\varepsilon=\sqrt\varepsilon$ and its added energy is
$O(r_\varepsilon^4)=o(r_\varepsilon^3)$.  It is compactly homotopic to the
identity, preserves degree, does not change the tangent defect measure, and
remains within the normalized additive quasiminimality error.

Thus GTA is false as a universal consequence of the AP-E17
\emph{additive} minimizing-diagonal axioms.  This does not disprove the
existence of a specially selected GTA recovery sequence or regularity of an
exact local minimizer.  The correct replacement gate is now explicit:
construct a scale-polished recovery sequence with a comparison gauge that is
uniform down to the needle scale, or prove the corresponding
energy--amplitude tightness directly.  All downstream promotions remain
closed.
\end{abstract}

\tableofcontents

\section{The exact logical input from AP-E17}

The frozen complete-minor density is
\begin{equation}
 W(F)=\frac12\left(|F|^2+R|M_2(F)|^2+K|M_3(F)|^2\right),
 \qquad R=1,\quad K=0.35.                              \label{eq:W}
\end{equation}
Let $u_{j_\ell}$ be the AP-E17 diagonal at a volumetric tangent point
$x_0$, with tangent radii $r_\ell\downarrow0$.  Its local relative-homotopy
deficit satisfies, for every fixed $0<\rho<1$,
\begin{equation}
 0\leq\epsilon_{j_\ell}(B_{\rho r_\ell}(x_0))
 \leq \cE(u_{j_\ell})-I_B
 =o\!\left(\mu(B_{r_\ell}(x_0))\right)
 =o(r_\ell^3).                                        \label{eq:additive}
\end{equation}
This is a strong statement at the tangent-ball scale.  It is only an
\emph{additive} statement: if $s_\ell\ll r_\ell$, then
\eqref{eq:additive} gives no useful estimate for
$\epsilon_{j_\ell}(B_{s_\ell})/s_\ell^3$.

In scaled coordinates put
\begin{equation}
 v_\ell(y)=\frac{u_{j_\ell}(x_0+r_\ell y)-q}{r_\ell},
 \qquad v_0(y)=F_0y,\qquad z_\ell=v_\ell-v_0.           \label{eq:tangent}
\end{equation}
The GTA weighted term is
\begin{equation}
 \cA_\ell(\delta)=
 \frac1{\delta^2}\int_{A_{2\delta}}
 |z_\ell|^2\left(1+R|Dv_\ell|^2
                 +K|M_2(Dv_\ell)|^2\right)\dd y .      \label{eq:GTA}
\end{equation}
AP-E17 proves the full boundary transfer if there exists
$\delta_\ell\downarrow0$ for which \eqref{eq:GTA} tends to zero, the
annular graph energy vanishes, and the annulus stays in one target
hemisphere.

\section{Why a standard reverse-H\"older theorem cannot be imported}

\subsection{The ambient \texorpdfstring{\((p,q)\)}{(p,q)} exponent test}

As a function of the ordinary derivative,
\begin{equation}
 c|F|^2\leq 1+W(F)\leq C(1+|F|^6).                     \label{eq:pq}
\end{equation}
Thus the coarse ambient exponents are $(p,q)=(2,6)$.  The 2024 partial
regularity theorem for relaxed strongly quasiconvex $(p,q)$ functionals
requires
\begin{equation}
 1\leq p\leq q<
 \min\left\{\frac{np}{n-1},p+1\right\}.                \label{eq:GK-range}
\end{equation}
For $n=3$ and $p=2$ the right side is $3$, whereas $q=6$.  Moreover, that
theorem explicitly does not derive the higher-gradient integrability which
would be needed to bridge this gap \cite{GmeinederKristensen2024}.

\subsection{The tempting \texorpdfstring{\(\mathscr A\)}{A}-free lift and its obstruction}

The complete-minor vector
\begin{equation}
 \Xi(Du)=\bigl(Du,\sqrt R\,M_2(Du),\sqrt K\,M_3(Du)\bigr) \label{eq:Xi}
\end{equation}
does satisfy linear differential identities: the first block is curl free,
and the Hodge duals of the second-minor forms are divergence free.  Since
$W=|\Xi|^2/2$, the 2025 $\mathscr A$-free higher-integrability theorem is a
natural candidate \cite{Schiffer2025}.

It does not apply to the present minimizer.  The range of $\Xi$ is not the
full kernel of a linear constant-rank operator.  It also satisfies
\begin{equation}
 M_2=Du\wedge Du,\qquad
 M_3=Du\wedge Du\wedge Du,\qquad
 |u|=1,\qquad u^TDu=0,                                 \label{eq:nonlinear}
\end{equation}
which are nonlinear algebraic and target constraints.  Local minimality is
known only among fields obeying \eqref{eq:nonlinear}, not among all
$\mathscr A$-free variations.  Enlarging the competitor class would assume
the very graph-recovery theorem under investigation.

\begin{proposition}[Reverse-H\"older import verdict]
Neither \eqref{eq:GK-range} nor the quadratic $\mathscr A$-free theorem
supplies a reverse-H\"older estimate for the AP-E17 minimizing diagonal.
Any valid proof must use new scale-aware minimality on the nonlinear
complete-minor graph.
\end{proposition}

\section{The boundary-straddling rank-two needle}

Fix a point $y_\ast\in\partial B_1$.  In a boundary chart, choose
$\psi,\phi\in C_c^\infty(B_1)$ such that on a smaller open set crossing the
flat interface through the origin,
\begin{equation}
 |\psi|+|D\psi|\geq c,\qquad |D\psi\wedge D\phi|\geq c. \label{eq:profiles}
\end{equation}
Take two orthonormal target directions $e_1,e_2\in T_q\SSS^3$.  For
$\varepsilon\downarrow0$ set
\begin{equation}
 A_\varepsilon=\varepsilon^{-1/2},\qquad
 b_\varepsilon=\varepsilon^2,\qquad
 h_\varepsilon=\varepsilon^2,                         \label{eq:scales}
\end{equation}
and define the tangent perturbation
\begin{equation}
 z_\varepsilon(y)=
 A_\varepsilon\psi\!\left(\frac{y-y_\ast}{h_\varepsilon}\right)e_1+
 b_\varepsilon\phi\!\left(\frac{y-y_\ast}{h_\varepsilon}\right)e_2 .
 \label{eq:needle}
\end{equation}
The support is a ball of radius $h_\varepsilon$ straddling
$\partial B_1$.  On the set in \eqref{eq:profiles}, $Dz_\varepsilon$ has
rank exactly two.

\begin{theorem}[Exact needle scaling]
Up to profile-dependent positive constants,
\begin{align}
 \int|z_\varepsilon|^2
 &\asymp A_\varepsilon^2h_\varepsilon^3
 =\varepsilon^5,                                      \label{eq:L2}\\
 \int|Dz_\varepsilon|^2
 &\asymp A_\varepsilon^2h_\varepsilon
 =\varepsilon,                                        \label{eq:E1}\\
 \int|M_2(Dz_\varepsilon)|^2
 &\asymp
 \frac{A_\varepsilon^2b_\varepsilon^2}{h_\varepsilon}
 =\varepsilon,                                        \label{eq:E2}\\
 M_3(Dz_\varepsilon)&=0.                              \label{eq:E3}
\end{align}
Nevertheless,
\begin{align}
 \int|z_\varepsilon|^2|Dz_\varepsilon|^2
 &\asymp A_\varepsilon^4h_\varepsilon=1,               \label{eq:weighted1}\\
 \int|z_\varepsilon|^2|M_2(Dz_\varepsilon)|^2
 &\asymp
 A_\varepsilon^2
 \frac{A_\varepsilon^2b_\varepsilon^2}{h_\varepsilon}
 =1.                                                   \label{eq:weighted2}
\end{align}
\end{theorem}

\begin{proof}
The support volume is comparable to $h_\varepsilon^3$.  The two derivative
sizes are $A_\varepsilon/h_\varepsilon$ and
$b_\varepsilon/h_\varepsilon=1$.  Equations
\eqref{eq:L2}--\eqref{eq:E2} follow by multiplying the corresponding
pointwise powers by the support volume.  The derivative has only two target
rows, so every pure third minor is zero.  On the active profile set,
$|z_\varepsilon|\asymp A_\varepsilon$; multiplying
\eqref{eq:E1} and \eqref{eq:E2} by $A_\varepsilon^2$ gives
\eqref{eq:weighted1}--\eqref{eq:weighted2}.
\end{proof}

Adding a bounded affine matrix $F_0$ does not alter these powers.  Since the
rank-two update has no term cubic in itself,
\begin{align}
 M_2(F_0+Dz)-M_2(F_0)&=L_2(F_0,Dz)+M_2(Dz),\notag\\
 M_3(F_0+Dz)-M_3(F_0)&=L_3(F_0,F_0,Dz)
 +L_3(F_0,Dz,Dz),                                    \label{eq:minor-expand}
\end{align}
and every displayed contribution has squared integral $O(\varepsilon)$.

\section{Failure for every shrinking annulus}

Let $0<\delta_\varepsilon<1/4$.  The active region in the intersection of
the needle and $A_{2\delta_\varepsilon}$ has volume bounded below by
\[
 c h_\varepsilon^3\quad\hbox{if }\delta_\varepsilon\geq h_\varepsilon,
 \qquad
 c h_\varepsilon^2\delta_\varepsilon
 \quad\hbox{if }\delta_\varepsilon<h_\varepsilon .
\]
It follows from \eqref{eq:weighted1} that
\begin{equation}
 \cA_\varepsilon(\delta_\varepsilon)\geq
 \begin{cases}
 c\,\delta_\varepsilon^{-2},
   &\delta_\varepsilon\geq h_\varepsilon,\\[1mm]
 c\,(\varepsilon^2\delta_\varepsilon)^{-1},
   &0<\delta_\varepsilon<h_\varepsilon.
 \end{cases}                                           \label{eq:all-delta}
\end{equation}
Both branches diverge for every choice
$\delta_\varepsilon\downarrow0$.  In contrast, the unweighted annular graph
energy obeys
\begin{equation}
 \int_{A_{2\delta_\varepsilon}}W(F_0+Dz_\varepsilon)
 \leq C
 \begin{cases}
 \varepsilon,&\delta_\varepsilon\geq h_\varepsilon,\\
 \delta_\varepsilon/\varepsilon,
    &\delta_\varepsilon<h_\varepsilon,
 \end{cases}
 \longrightarrow0,                                    \label{eq:ann-energy}
\end{equation}
because the second branch has
$\delta_\varepsilon<h_\varepsilon=\varepsilon^2$.

\begin{theorem}[GTA instability]
The strong tangent convergence, vanishing annular graph energy, and
vanishing normalized added energy do not imply GTA, even for a genuinely
rank-two perturbation whose physical image stays in one target hemisphere.
\end{theorem}

\section{Compatibility with an additive minimizing diagonal}

Take the physical tangent radius to be
\begin{equation}
 r_\varepsilon=\varepsilon.                            \label{eq:r}
\end{equation}
First insert \eqref{eq:needle} through the exponential chart over the
bounded affine background $v_0=F_0y$.  More generally the same calculation
holds on a recovery patch whose background complete-minor graph is
equiintegrable at the needle scale.  Its largest target displacement is
\begin{equation}
 r_\varepsilon A_\varepsilon=\sqrt\varepsilon\to0,     \label{eq:hemisphere}
\end{equation}
so radial projection and its exterior powers have uniform constants.  The
physical added derivative energy is
\begin{equation}
 r_\varepsilon^3\,O(\varepsilon)
 =O(r_\varepsilon^4)=o(r_\varepsilon^3).               \label{eq:physical}
\end{equation}
The perturbation is supported in a contractible target chart and the
homotopy obtained by multiplying \eqref{eq:needle} by $t\in[0,1]$ is compact.
It therefore preserves the exterior trace and global degree.

\begin{theorem}[Instability of the AP-E17 additive tangent axioms]
On the affine tangent background, the rank-two needle has
$o(r_\varepsilon^3)$ physical cost.  It does not change the weak tangent
graph measure, and its removable energy is compatible with
\begin{equation}
 \frac{\epsilon_{j_\ell}(B_{\rho r_\ell})}
      {\mu(B_{r_\ell})}\longrightarrow0
\end{equation}
for every fixed $\rho$ whenever it is inserted into a controlled recovery
patch.  Nevertheless it violates
\eqref{eq:GTA} for every $\delta_\ell\downarrow0$.
\end{theorem}

\begin{proof}
For the affine background, \eqref{eq:minor-expand} and
\eqref{eq:E1}--\eqref{eq:E2} make the normalized graph-energy change
$O(\varepsilon)$; hence it disappears from the weak tangent measure and is
$o(1)$ at the additive comparison normalization.  The same conclusion holds
on a patch for which the background mixed terms in
\eqref{eq:minor-expand} are equiintegrable at that scale.  Degree and trace
are unchanged by the compact homotopy.  Equation \eqref{eq:all-delta} proves
failure of the remaining GTA condition.
\end{proof}

This theorem has a deliberately precise scope.  It disproves a theorem
which tries to derive GTA from the AP-E17 additive tangent estimates alone:
those estimates cannot distinguish the needle.  It does not prove that the
repository's selected diagonal actually contains the needle, that no
carefully selected recovery sequence has GTA, or that the degree-one
minimizer has a nonzero relaxation defect.

\section{The corresponding reverse-H\"older no-go}

On its active set the needle graph density satisfies
\begin{equation}
 G_\varepsilon\asymp\varepsilon^{-5},
 \qquad
 |\operatorname{supp}G_\varepsilon|\asymp\varepsilon^6,
 \qquad
 \int G_\varepsilon\asymp\varepsilon.                  \label{eq:G}
\end{equation}
For every $\sigma>0$,
\begin{equation}
 \int G_\varepsilon^{1+\sigma}
 \asymp\varepsilon^{1-5\sigma},\qquad
 \frac{\left(\int G_\varepsilon^{1+\sigma}\right)^{1/(1+\sigma)}}
      {\int G_\varepsilon}
 \asymp\varepsilon^{-6\sigma/(1+\sigma)}
 \longrightarrow\infty.                               \label{eq:RHfail}
\end{equation}
Thus no sequence-uniform fixed-ball $L^{1+\sigma}$ gain follows from
\eqref{eq:additive}.  At the needle scale the removable energy is
$O(\varepsilon)$, exactly the size allowed by the additive error.

\section{The scale-polished repair}

\begin{definition}[Uniform microscale comparison gauge]
A tangent recovery sequence is scale polished if there are
$\omega_\ell\downarrow0$ and $s_\ell\downarrow0$ such that for every ball
$B_s(y)\Subset B_1$ with $s_\ell\leq s\leq1/4$,
\begin{equation}
 \epsilon_\ell(B_s(y))
 \leq \omega_\ell s^3.                                \label{eq:polished}
\end{equation}
An exact local minimizer has $\omega_\ell=0$.
\end{definition}

The needle violates \eqref{eq:polished} at
$s=h_\varepsilon=\varepsilon^2$, since its removable energy is
$\asymp\varepsilon$ whereas $s^3=\varepsilon^6$.  Thus
\eqref{eq:polished} detects exactly what the global additive estimate misses.
It is not yet proved that a degree-preserving recovery sequence satisfying
\eqref{eq:polished} exists.

A second, weaker sufficient target is the energy--amplitude tightness
\begin{equation}
 \int_{B_1}|z_\ell|^2
 \left(1+R|Dv_\ell|^2+K|M_2(Dv_\ell)|^2\right)\dd y
 \longrightarrow0.                                    \label{eq:EAT}
\end{equation}
Together with no boundary mass and a one-hemisphere collar,
\eqref{eq:EAT} permits a slow diagonal choice
$\delta_\ell\downarrow0$ with
\[
 \delta_\ell^{-2}\times\text{the left side of \eqref{eq:EAT}}\to0
\]
and vanishing annular graph energy, hence GTA.  The needle has left side
$\asymp1$ and is the sharp obstruction.

\section{Gate decision}

The requested GTA investigation therefore closes with a negative universal
result and a narrower positive target:
\begin{enumerate}
 \item \textbf{Reverse-H\"older from the current diagonal: false.}
 The additive error is not scale uniform; current $(p,q)$ and
 $\mathscr A$-free theorems do not apply.
 \item \textbf{GTA from the additive tangent axioms alone: false.}
 The boundary-straddling rank-two needle is invisible to those normalized
 limits while violating every annulus thickness.  Whether the selected
 repository diagonal contains or can be polished to exclude it remains open.
 \item \textbf{GTA for a scale-polished diagonal: open.}
 The next theorem must construct \eqref{eq:polished}, or prove
 \eqref{eq:EAT} plus the hemisphere collar by another exact local-selection
 argument.
\end{enumerate}

The counterexample does not reopen lattice work and does not authorize the
bosonic Hessian.  Full $\mu=0$, local recovery, classicality, continuum
isolation, determinant promotion, and the degree-one portal remain false.
Parallel Dirac/Callias mathematics on a prescribed smooth background remains
allowed.

The independent verifier
\begin{center}
\texttt{route\_f/code/verify\_ap\_e18\_gta\_minimizing\_diagonal.py}
\end{center}
checks the rank, every scaling exponent, both annulus branches,
reverse-H\"older divergence, literature applicability, and all fail-closed
gates.

\begin{thebibliography}{9}

\bibitem{GmeinederKristensen2024}
F.~Gmeineder and J.~Kristensen,
\newblock Quasiconvex functionals of \((p,q)\)-growth and the partial
regularity of relaxed minimizers,
\newblock \emph{Arch. Rational Mech. Anal.} \textbf{248} (2024), 80,
\newblock \url{https://doi.org/10.1007/s00205-024-02013-8}.

\bibitem{Schiffer2025}
S.~Schiffer,
\newblock \(\mathscr A\)-free truncation and higher integrability of
minimisers,
\newblock \emph{Calc. Var. Partial Differential Equations} \textbf{64}
(2025), 30,
\newblock \url{https://doi.org/10.1007/s00526-024-02855-w}.

\end{thebibliography}

\end{document}
